\chap{UI 与汇报生成指南：从共同原则到典型表达}
\label{sec:guidelines}
\subsection{六项共同原则}
第一，围绕人的理解和决策组织内容。生成前明确人需要知道什么、为什么需要知道，以及理解后可以做什么。首屏应建立足够背景并突出当前重点。对于架构讨论，重点可能是职责和依赖；对于验收，重点可能是可体验的变化和待判断的取舍。

第二，让理解和控制共处。视图中的对象应支持适合当前任务的操作，例如体验、比较、标注和选择。反馈携带必要上下文返回 AI。若某视图只用于讲解，也应提供继续提问或定位相关工作的入口，使解释能够支持后续行动。

第三，为信息提供背景和位置。设计点关联它涉及的模块、行为和目标，改进说明当前与拟议状态的差异。简短标签可以帮助导航，关键机制则需要实例和关系解释。对于熟悉背景的人可以压缩呈现，对首次接触者应允许展开基本概念。

第四，支持逐步探索并保持方向感。人可以从整体进入局部，从结果进入机制，从结论进入依据。深入后仍能知道当前位置和与整体的关系。跨视图跳转时保留当前对象、场景与筛选条件，减少重复定位。

第五，让表达与实际工作对应。当前实现、拟议方案、模拟结果和已验证事实具有不同状态，操作也需要明确作用范围。证据入口帮助人核对重要结论，状态变化则应来自实际执行结果。准确性是可理解性的基础。

第六，根据使用反馈改进表达。反复追问、错误理解、难以找到操作和中断后的恢复困难，都可以作为改进线索。常用概念与入口保持连贯，新的表达围绕具体问题展开。评价改进时观察理解与操作效果，而不单独追求新颖性。

\subsection{指南一：层次化架构与协议探索}
这类视图适用于理解系统组成、判断职责边界和评估改动影响。初次进入时显示主要功能或服务模块及其责任，避免同时展开全部内部细节。进入一个模块后，显示内部组成、关键数据和对外联系，并保留整体位置。

模块之间的边可以成为协议入口。点击之后，用一个请求如何经过双方、产生什么结果来解释契约，进一步展示异常和并发分支。协议字段可以按需展开，先让人理解双方各自保证什么。改进建议附着在受影响的节点和关系上，使变化具有可见的位置。

\fig{10-architecture}{架构理解指南。左侧逐级呈现模块职责，右侧以具体交互解释选中的协议；拟议改进关联到它影响的对象。}{architecture}

这类视图需要同时保留层次和跨模块关系。目录树适合定位内部组成，网络关系适合追查跨边界影响。两者可以通过按需展开结合。若人无法解释某个模块为何独立，或频繁需要穿越多个模块理解一个内部决定，视图可以帮助发现值得讨论的边界问题。

\subsection{指南二：围绕真实任务的功能体验}
功能体验视图围绕一个明确任务组织操作、反馈和结果。例如，让读者启动一个模拟任务，提出取消，观察等待和最终状态，再查看保留输出。用户通过实际尝试建立理解，说明文字则解释关键差异和限制。

对比改进前后行为时，应使用相同任务条件。两个版本的不同输入会使人难以判断差异来自哪里。可以提供情境切换，例如执行器延迟、重复点击或任务已经完成。反馈入口附着在具体步骤上，使人能够指出“此时我以为已经停止”之类的体验问题。

\fig{11-workbench}{可操作的功能体验示意。任务状态、演示操作、机制说明和针对步骤的反馈共处；图中的控件属于模拟演示设计。}{workbench}

评价这类视图时，观察读者能否独立完成任务、是否正确理解结果，以及反馈是否能够定位到可修改对象。演示如果过于理想化，可能隐藏真正困难的分支；因此需要让重要异常情境可以被尝试。

\subsection{指南三：运行机制与可调情境}
机制视图适用于解释状态、事件顺序、资源变化和失败处理。它可以采用时序图、状态机或逐步动画，并允许暂停、回退和调整条件。一个状态节点应说明进入条件和退出条件，事件边则说明由谁发出、携带什么含义。

机制解释的关键是将抽象规则放进具体过程。例如，取消请求与完成事件竞争时，改变二者到达顺序，观察终态如何确定。人由此理解系统为何需要区分接收与停止，而不必先阅读完整接口文档。模拟过程应明确其假设，实际运行证据则作为另一类入口。

人可以在异常分支上提出新的约束，例如资源释放失败时需要显示什么信息。这样的反馈直接作用于机制设计，并能进入相应测试任务。若动画只是沿箭头播放而无法解释条件，它的视觉变化并没有增加多少理解价值。

\subsection{指南四：共同场景下的方案比较}
方案比较视图适用于需要人表达偏好或承担关键取舍的情境。候选方案应面对相同目标与场景，呈现行为、收益、代价和重要不确定性。对无法测量的效果应使用定性解释，并说明依赖条件，避免用精确数字制造确定感。

例如，取消后保留全部中间结果与只保留可解释部分，会影响存储成本、结果质量和用户预期。界面可以让人查看同一任务在两种方案下的结果组织，再选择偏好或补充约束。选择结果应关联其理由，便于后来在条件变化时重新评估。

比较视图不必要求人决定每个实现细节。对已被既定偏好覆盖的重复取舍，Agent 可以自行处理并在需要时解释；对影响目标或资源承诺的差异，再提供聚焦的选择界面。其价值在于提高判断质量并减少无关决策。

\subsection{指南五：阶段进展与注意力分配}
进展视图可以围绕人上次查看之后的重要变化展开，呈现可体验成果、仍然存在的阻塞和需要关注的取舍。完成状态应关联实际交付与验证范围，避免把“正在准备”“已经提交修改”和“用户可用”混为一谈。

人的关注点不同，进展选择也应不同。负责产品方向的人可能关心行为变化和使用反馈，关注架构的人可能关心职责与协议变化。系统可以提供默认重点和完整入口，让人按需展开。重点选择应能解释其理由，而不是隐藏所有不利结果。

操作可以包括体验成果、追查阻塞、调整优先级和回应问题。需要人回应的事项应足够具体，提供背景及可比较选项。大量“请确认”卡片会重新形成控制负担，因此日常工程判断应尽量在已授权范围内完成。

\subsection{指南六：未来演化与依赖路径}
未来演化视图在当前系统上叠加拟议变化，展示它们服务的目标、前置条件和可能影响。人可以切换不同路径，观察先做哪项工作会支持或限制后续改进。拟议节点应与当前实现具有清楚的视觉状态，避免让读者误以为能力已经存在。

这种视图可以把零散改进点组织成可讨论的路径。例如，先明确取消协议，再改善状态投影，随后增加可解释的输出保留策略。若某项基础工作能支持多个目标，其价值可以通过依赖关系解释；若某个方案需要大量前置修改，也应让人看见代价。

\fig{12-guides}{六类典型投影按理解任务组织。每类同时给出探索过程与人的参与方式；实际 Workbench 可以组合多个指南，并根据反馈改写。}{guides}

\subsection{把指南转化为可执行的生成任务}
以架构指南为例，生成任务可以给出目标读者、当前疑问、相关对象、已知事实、可用动作和希望核对的理解。假设读者熟悉用户功能，却不了解执行器协议，任务重点就是让其理解“取消请求为何不等于已经停止”。所需输入是有关模块及协议、典型事件序列和一个延迟停止情境。

Agent 可以先生成模块层次与协议入口，再提供可暂停的事件演示，最后把反馈动作绑定到选中的状态。核对时请读者判断某一时刻任务是否已停止，并说明依据。如果读者只记住图上的颜色，却不能理解状态条件，就应修改解释和交互。任务以需要建立的理解约束结果，具体视觉实现仍保持开放。

同样的任务可以面向不同场景改写。阶段汇报可采用较紧凑的静态图与关键片段，深入讨论可提供可调参数的演示，移动设备上则可以按步骤展开。生成协议只需保持对象、状态与动作的一致含义，避免将某种布局固化为所有情境的唯一模板。

\subsection{指南的组合、生成与修订}
每份指南可以使用“适用疑问、首先呈现什么、怎样深入、人怎样参与、怎样评价”五项说明。它约束需要建立的理解，给具体布局和媒介组合留下空间。模板可以固化常用对象和交互结构，skill 可以说明生成与核对方法，协议可以规定状态与动作的语义联系。

生成时，Agent 先确定当前问题与事实依据，再选择适当指南，构造可查看的初版，检查关键文字和动作，最后根据人类反馈调整。若事实不足，可以生成标明假设的方案演示；若对象状态已经变化，则应更新相应部分。生成和核对本身也可以由不同成员协作完成。

跨视图组合应围绕人的疑问推进。例如，从阶段进展中的延迟问题进入机制演示，发现等待关系后查看架构协议，再比较两个修改方案。对象与场景随跳转保留，能够减少重复解释。有效形式无法穷举，指南应作为经验的积累方式，持续接受新的表达与反例。
